% Part: first-order-logic
% Chapter: axiomatic-deduction
% Section: axioms-rules-propositional

\documentclass[../../../include/open-logic-section]{subfiles}

\begin{document}

\iftag{FOL}
      {\olfileid{fol}{axd}{prp}}
      {\olfileid{pl}{axd}{prp}}

\olsection{વિધાનાત્મક સંયોજકો માટે સ્વયંસિદ્ધિઓ અને નિયમો}

\begin{defn}[સ્વયંસિદ્ધિઓ]
વિધાનાત્મક સંયોજકો માટેની \emph{સ્વયંસિદ્ધિઓ}નો ગણ~$\PAx$ નીચેના
સ્વરૂપોનાં બધાં !!{formula}sથી બને છે:
\begin{align}
  & (!A \land !B) \lif !A \ollabel{ax:land1}\\
  & (!A \land !B) \lif !B \ollabel{ax:land2}\\
  & !A \lif (!B \lif (!A \land !B)) \ollabel{ax:land3}\\
  & !A \lif (!A \lor !B) \ollabel{ax:lor1}\\
  & !A \lif (!B \lor !A) \ollabel{ax:lor2}\\
  & (!A \lif !C) \lif ((!B \lif !C) \lif ((!A \lor !B) \lif !C)) \ollabel{ax:lor3}\\
  & !A \lif (!B \lif !A) \ollabel{ax:lif1}\\
  & (!A \lif (!B \lif !C)) \lif ((!A \lif !B) \lif (!A \lif !C)) \ollabel{ax:lif2}\\
  & (!A \lif !B) \lif ((!A \lif \lnot !B) \lif \lnot !A) \ollabel{ax:lnot1}\\
  & \lnot !A \lif (!A \lif !B) \ollabel{ax:lnot2}\\
  & \ltrue \ollabel{ax:ltrue}\\
  & \lfalse \lif !A \ollabel{ax:lfalse1}\\
  & (!A \lif \lfalse) \lif \lnot !A \ollabel{ax:lfalse2}\\
  & \lnot\lnot !A \lif !A \ollabel{ax:dne}
\end{align}
\end{defn}

\sourcecorrection{OLAX-001}{મૂળની
\texttt{axioms-rules-propositional.tex}, પંક્તિ~16માં
``The set of $\PAx$ of axioms'' લખતાં વધારાનું ``of'' આવ્યું છે. અહીં
$\PAx$ને વિધાનાત્મક સંયોજકો માટેની સ્વયંસિદ્ધિઓના ગણ તરીકે સ્પષ્ટ
લખ્યું છે.}

\begin{defn}[મોડસ પોનેન્સ]
  જો $!B$ અને $!B \lif !A$ પહેલેથી !!a{derivation}માં આવતાં હોય, તો
  $!A$ યોગ્ય અનુમાન-પગલું છે.
\end{defn}

મોડસ પોનેન્સના નિયમને આપણે ટૂંકમાં ``\MP'' લખીશું.

\end{document}
